\section{Atak bruteforce}

Atak odbył się przy pomocy prostego skryptu sprawdzającego, czy kolejne hasła z listy 1050 popularnych dają dostęp do wybranego konta.
\begin{table}[H]
\begin{tabular}{|l|l|l|l|l|l|}
\hline
\rowcolor[HTML]{EFEFEF} 
\multicolumn{1}{|c|}{\cellcolor[HTML]{EFEFEF}\textbf{Metoda}} & \multicolumn{1}{c|}{\cellcolor[HTML]{EFEFEF}\textbf{Końcówka}} & \multicolumn{1}{c|}{\cellcolor[HTML]{EFEFEF}\textbf{Nagłowki}} & \multicolumn{1}{c|}{\cellcolor[HTML]{EFEFEF}\textbf{Ciasteczka}} & \multicolumn{1}{c|}{\cellcolor[HTML]{EFEFEF}\textbf{Dane}}                             & \multicolumn{1}{c|}{\cellcolor[HTML]{EFEFEF}\textbf{Odpowiedź/wynik}} \\ \hline
GET    & 
/rest/user/login  & -   & -   & \begin{tabular}[c]{@{}l@{}}\{'email':admin@juice-sh.op, 'password':admin123\}\end{tabular} &  \begin{tabular}[c]{@{}l@{}}w tym przykładzie kod 200, bo hasło poprawne \\ w pozostałych kod 401 z błędem \end{tabular} \\ \hline
\end{tabular}
\end{table}

Łamanie hasła tą metodą w wybranym uruchomieniu zajęło 116 prób w  9.05~s, co daje ~0.078s na próbę - znacznie szybciej niż mógłby to wykonać człowiek, czyli aplikacja jest podatna na bruteforce.


{\rule{\linewidth}{0.2mm}}
\begin{center}
    \textcolor{red}{APLIKACJA PODATNA NA BRUTEFORCE LOGIN}
\end{center} 